\subsection{The inherited cost of arithmetic slats}
\label{geomtransfer:slats}

The sharp finite train-track estimate uses a global arithmetic slat
average. The following obstruction explains why that average cannot be
inserted unchanged at every stage of one fixed singular construction.

\begin{proposition}[No uniform full-grid slats for a fixed measure]
\label{geomtransfer:slatrigidity}
Let \(\sigma\) be a finite measure on a compact interval, and let
\(h_j\to0\), \(\delta_j/h_j\to0\). Let \(U_j\) be the union of
at most \(C_0/h_j\) intervals \(I_{j,k}\), each of length at most
\(C_0\delta_j\), whose distinct centers belong to an arithmetic grid of
spacing \(h_j\). There do not exist positive submeasures
\(\sigma_j\le\sigma\), supported on \(U_j\), for which
\[
 \sigma_j(\mathbb R)\ge m_0>0,
 \qquad \sigma_j(I_{j,k})\le C_1h_j
\]
with constants independent of \(j\).
In particular a fixed probability cannot itself satisfy these global
slat hypotheses at infinitely many resolutions.
\end{proposition}

\begin{proof}
For large \(j\) the slats are disjoint, and any interval \(J\)
meets at most \(C(|J|/h_j+1)\) of them. Thus
\[
 \sigma_j(J)\le CC_1(|J|+h_j).
\]
Pass to a weakly convergent subsequence on the common compact support.
The limit \(\tau\) has mass at least \(m_0\). Applying the preceding
bound to open intervals and using weak convergence gives
\(\tau(J)\le CC_1|J|\), so \(\tau\) has a bounded Lebesgue density.

Write \(\sigma=\sigma^{\mathrm{ac}}+\sigma^{\mathrm{s}}\).
The slat unions have length at most
\(C_0^2\delta_j/h_j\to0\). Hence the absolutely continuous part
of \(\sigma_j\) has mass at most
\(\sigma^{\mathrm{ac}}(U_j)\to0\). Its singular part is dominated
by \(\sigma^{\mathrm{s}}\). Testing nonnegative continuous functions
and passing to the limit therefore gives \(\tau\le\sigma^{\mathrm{s}}\).
This contradicts the positive mass and absolute continuity of \(\tau\).
\end{proof}

The inherited normalization loss is quantitative. Suppose the old vertical
support lies in \(J\) intervals of total length \(L\), and a retained
measure \(\tau\) of mass at least \(m_0\) satisfies
\[
 \tau(I)\le C(|I|+h)\qquad\text{for every interval }I.
\]
Summing over the old parent intervals forces
\begin{equation}\label{geomtransfer:inheritedcost}
 C\ge\frac{m_0}{L+Jh}.
\end{equation}
In an interval Moran construction with \(M_i\) children of length
\(\ell_i\) and spacing \(h_i=\ell_{i-1}/M_i\), the occupied length
after stage \(j\) is
\[
 L_j=\ell_0\prod_{i\le j}\frac{\ell_i}{h_i}.
\]
If \(L_j\to0\), later full-grid counting constants must grow at least
as a multiple of \(L_j^{-1}\), once the fresh mesh error is negligible.
Conditioning on an old interval of length \(\ell\) instead gives a
normalized count proportional to \((|I|+h)/\ell\); that parent factor
cannot be dropped.

This does not prove that actual collision phase windows attain the worst
interval in \eqref{geomtransfer:inheritedcost}, or that the loss survives
after all conditional source and pin weights are summed. It rules out
literal reuse of the uniform finite average with the same constant.
Controlling those weighted inherited costs remains necessary for a
nontrivial nested transfer.

\subsection{A quantitative pinned estimate for generic nearby metrics}
\label{geomtransfer:metrics}

Positivity of distance sets for almost every elliptic metric is established
context \cite{hofmanniosevich}; no novelty for that phenomenon is claimed
here. The following complete estimate shows why a fixed profile witness
is not itself a Euclidean null-distance example.

Let \(\mu,\nu\) be compactly supported planar probabilities, with
\(I_1(\mu)<\infty\). The pin probability is fixed independently of the
metric parameter and may equal \(\mu\). For \(0<\eta<1/2\), let
\(\mathbb P_\eta\) be normalized Lebesgue measure on the Frobenius ball
\[
 \mathcal M_\eta=\{M\in\operatorname{Sym}_2(\mathbb R):
                          \|M-I\|_F<\eta\}.
\]
Its matrices are positive definite. Put
\(d_{M,y}(x)=((x-y)^{\top}M(x-y))^{1/2}\).

\begin{proposition}[Mean joint pinned second norm]
\label{geomtransfer:metricbound}
The joint distance law has a density with respect to
\(d\mathbb P_\eta(M)\,d\nu(y)\,dt\), and
\begin{equation}\label{geomtransfer:metricnorm}
 \int_{\mathcal M_\eta}\int
       \|(d_{M,y})_*\mu\|_2^2\,d\nu(y)\,d\mathbb P_\eta(M)
 \le C\eta^{-1}I_1(\mu).
\end{equation}
No source--pin separation is required. In particular, for almost every
matrix, the full pinned law has an \(L^2\) density for
\(\nu\)-almost every pin.
\end{proposition}

\begin{proof}
For a source pair and pin, write
\[
 v=x-x',\qquad z=(x+x')/2-y,
 \qquad Q=zv^{\top}+vz^{\top}.
\]
The squared-distance difference is \(\operatorname{tr}(MQ)\), and
\begin{equation}\label{geomtransfer:matrixnorm}
 \|Q\|_F^2=2|z|^2|v|^2+2(z\cdot v)^2.
\end{equation}
Since all eigenvalues of \(M\) lie between \(1/2\) and \(3/2\),
\[
 d_{M,y}(x)+d_{M,y}(x')\le C(|v|+|z|).
\]
A distance difference at most \(\varepsilon\) therefore implies
\( |\operatorname{tr}(MQ)|\le C\varepsilon(|v|+|z|)\).
The proportion of a three-dimensional ball of radius \(\eta\)
lying in a linear slab of width \(a\) is at most \(Ca/\eta\).
Together with \eqref{geomtransfer:matrixnorm}, this yields
\begin{equation}\label{geomtransfer:slabbound}
 \varepsilon^{-1}\mathbb P_\eta
 \{|d_{M,y}(x)-d_{M,y}(x')|\le\varepsilon\}
 \le C\eta^{-1}(|v|^{-1}+|z|^{-1})
\end{equation}
when neither vector vanishes.

The midpoint term cannot be discarded merely from separation of the two
supports. Instead, for reflection \(R_y(x)=2y-x\),
\begin{equation}\label{geomtransfer:midpoint}
 \iint |(x+x')/2-y|^{-1}\,d\mu(x)d\mu(x')
 =2I_1(\mu,(R_y)_*\mu)\le2I_1(\mu).
\end{equation}
The inequality is Cauchy--Schwarz for the positive Riesz-energy form:
its Fourier weight is a constant times \(|\xi|^{-1}\), and reflection
and translation preserve self-energy. Equivalently one may use the
positive Gaussian representation of the Riesz kernel. Finite energy
implies that \(\mu\) is nonatomic, so both the diagonal and, for each
fixed \(y\), the reflected diagonal have zero product mass. Thus the
exceptions in \eqref{geomtransfer:slabbound} are null after integration
against any prescribed \(\nu\).

Take a nonnegative smooth probability mollifier supported in
\([-1/2,1/2]\). Its self-correlation is bounded and supported in
\([-1,1]\), so the squared norm of a smoothed distance law is bounded
by a constant times its collision probability of width
\(\varepsilon\), divided by \(\varepsilon\). Integrating
\eqref{geomtransfer:slabbound} and using
\eqref{geomtransfer:midpoint} gives the uniform bound
\(C\eta^{-1}I_1(\mu)\) for these norms in the fixed product space
\(L^2(d\mathbb P_\eta\,d\nu\,dt)\).
The radial supports lie in a common bounded interval. Weak compactness,
and weak convergence of mollifications of the joint distance law,
identify its actual density. Lower semicontinuity gives
\eqref{geomtransfer:metricnorm}, and disintegration gives the
almost-everywhere assertions.
\end{proof}

Retaining the distance-sum factor in the slab calculation is essential;
replacing it by a diameter would introduce a product of reciprocal
lengths instead of their sum. For \(L_M=M^{1/2}\),
\[
 d_{M,y}(x)=|L_Mx-L_My|.
\]
Thus almost every matrix in any prescribed small ball produces one fixed
Euclidean image of the source and pin measures with the asserted raw
joint density. The linear map preserves their Hausdorff and packing
dimensions and their Frostman exponents, and differs from the identity
by \(O(\eta)\). This applies in particular to source measures with
dimension exponent greater than one, including the fixed realizability
constructions discussed earlier.

The norm bound deteriorates as \(\eta^{-1}\). Consequently taking good
matrices closer to the identity gives no uniform second-norm control and
does not establish the conclusion for the original Euclidean metric.
That metric may belong to the exceptional parameter set. This geometric
comparison supplies no additional dimension-only region for the original
set.
